\section{What Did Not Move Live Accuracy, and Why}\label{sec:v4negative}

This section reports the negative results of the V4 research wave that followed the registered look of \cref{sec:candidate}. Every
number is exploratory and, except for the ReLU$^2$ kill readings (trainer dev-clean WER), was read on the selection half of DEV (S) with the pinned Rust scorer and a paired speaker bootstrap of \VfBootReps{}
resamples; the unselected half (C) was used only for the diagnosis of the live gap, and the test sets were not read
(\cref{sec:protocol}). Differences are in percentage points of pooled WER at the live setting c4 (a 160\,ms chunk for RC2), negative meaning lower WER.

\para{Where the live gap sits.} Of the \VfDiagPositions{} C-half reference positions that the compact-LM live stream gets wrong and
Zipformer INT8 gets right, \VfDiagInterior{} (\VfDiagInteriorPct\%) are interior words, not the first or last word of an utterance
(\VfDiagEdge{} positions). \VfDiagFreqOther{} are dev-other interior substitutions of words that occur at least \VfFreqMin{} times in the training
transcripts, and of \VfDelTotal{} deletions only \VfDelLast{} are final words, so the \VfDelGap-deletion excess over Zipformer is
not a tail-drop. The residual therefore looks like acoustic discrimination of common words in hard audio, not
vocabulary, endpointing or boundaries; \VfCellCount{} such S-half positions were frozen as a secondary read.

\para{Reranking and LM cannot reach the target.} On S, the 1-best of the live stream with the compact LM makes
\VfOracleOneBestEdits{} edits (\VfOracleOneBestWer\% WER). A perfect ranker over its beam would reach \VfOracleDepthTwoWer\% at depth \VfDepthTwo{}
and \VfOracleDepthFourWer\% (\VfOracleDepthFourEdits{} edits) with the full beam of \VfDepthFour{}, against a
\VfOracleTargetWer\% target of \VfOracleTargetEdits{} edits. Reranking or a stronger live LM on this beam therefore cannot reach \VfOracleTargetWer\% WER (at most \VfOracleSecondaryEdits{}
edits of headroom). A width-1 fused greedy decoder with an LM was \VfNineDeltaEdits{} edits worse than greedy
decoding (\VfNineLmEdits{} vs.\ \VfNineGreedyEdits), missing its \VfNineBar\pp\ gain gate.

\para{Stronger teacher labels.} Continuing RC2 training on audio re-labelled by a stronger teacher (PL3) did not lower live WER
under the registered rule: the greedy c4 difference to the incumbent labels was \VfPlThreeDelta\pp\ (one-sided 95\% bounds \VfPlThreeLo{}
to \VfPlThreeHi\pp) against a required gain above \VfPlThreeBar\pp, an independently reproduced non-adoption.

\para{Architecture probes.} Converting the FFN activation to ReLU$^2$ from the RC2 weights, intended to create zero-skippable activations,
was stopped by its registered kill at optimizer step \VfReluKillStep: dev-clean EMA WER was \VfReluOneEma\% (raw \VfReluOneRaw\%) in one
arm, with \VfReluOneZero\% of the activations already zero, and raw WER was \VfReluTwoRaw\% in the other, far beyond the
\VfReluKillPp\pp\ allowance. Chunk-causal lookahead convolution ($k_\mathrm{right}{=}\VfBoneRight$) did not help either (below). In a small-model scout trained on LibriSpeech only, a TDT output head was
\VfTdtDelta\pp\ worse than exact RNN-T at that scout's own c4 setting (\VfTdtChunkMs\,ms; bounds \VfTdtLo{} to \VfTdtHi\pp; RNN-T reference \VfTdtRnntWer\% pooled, not comparable with RC2), so exact RNN-T stays the control.
Skipping the upper encoder on blank chunks after midpoint-CTC certification, with the refresh needed to keep the state exact,
netted only \VfSkipLo--\VfSkipHi\% skipped chunks against a \VfSkipGate\% gate, so it was dropped.

\para{Regularisers and the screen protocol.} Short fine-tuning screens (\VfScreenSteps{} steps from RC2, EMA weights) were each paired
against a matched control (S-NULL) with the same recipe, data and seed. A second control seed differed from the first by
\VfFloorDelta\pp\ (bounds \VfFloorLo{} to \VfFloorHi\pp), with seed and device confounded, so a single-seed screen cannot resolve effects
below about \VfSingleSeedPp\pp. Each arm reported here was therefore run with \VfScreenSeeds{} seeds. Only the chunk-causal convolution had a registered
bar, \VfScreenBar\pp; the other screens were descriptive and are read against the noise floor.
The chunk-causal convolution gave \VfBoneSeedA\pp{} and \VfBoneSeedB\pp{} on the two seeds (mean \VfBoneMean\pp) and
\VfBoneOnLd\pp{} on top of layer drop. Layer drop at rate \VfLdRate{} averaged \VfLdMean\pp, a SpecAugment time-mask budget capped at \VfSpecCapPct\% of frames
(\VfSpecMasks{} masks of up to \VfSpecWidth{} frames, \VfSpecMaskedPct\% of frames masked) \VfSpecMean\pp, and both together \VfComboMean\pp, of which adding the mask
cap to layer drop contributed \VfComboOnLd\pp. All of these lie inside the control noise floor and the chunk-causal convolution did not reach its bar; they are
no detectable effect, not a measured zero. The full-length layer-drop run (O3-LD, layer drop \VfLdRate{} over the full training schedule, seed \VfSeed) was read against
its registered rule on the same selection half: greedy c4 differed from its matched control by \VfOThreeDelta\pp (bounds
\VfOThreeLo{} to \VfOThreeHi\pp; \VfOThreeEdits{} vs.\ \VfOThreeControlEdits{} edits), against a required gain of
\VfOThreeBar\pp with an upper bound below zero; c8 and c16 (guard \VfOThreeGuard\pp) differed by \VfOThreeEightDelta\pp and
\VfOThreeSixteenDelta\pp. It was not adopted, and a second-seed replicate was stopped at the read, as registered (no result).

\para{Implication.} Within these screens, post-processing the beam, relabelling existing audio, lookahead convolution, and
regularisation did not move live accuracy, and the residual is concentrated in common interior words on hard audio. The remaining
levers are more acoustic evidence (fresh audio, not relabelled audio) and capacity paid for by the power savings of \cref{sec:versus}; we do not
yet have a result showing that either closes the gap.
